\documentclass[10pt, a4paper]{article}

% --- 1. 필수 패키지 설정 ---
\usepackage[a4paper, top=18mm, bottom=15mm, left=12mm, right=12mm, headheight=25pt, headsep=8mm]{geometry}
\usepackage{amsmath, amssymb, amsthm}
\usepackage{fancyhdr}
\usepackage{enumitem}
\usepackage{array}
\usepackage{longtable}
\usepackage{kotex}
\usepackage{setspace}

% --- 2. 헤더 설정 ---
\pagestyle{fancy}
\fancyhf{}
\lhead{\large \textbf{미적분 : 제15장 정적분의 계산 연습문제}}
\rhead{\small 학번: \underline{\hspace{3.5cm}} 이름: \underline{\hspace{3cm}}}
\cfoot{\small \thepage}
\renewcommand{\headrulewidth}{0.6pt}

% --- 3. 문제 박스 명령어 정의 ---
\newcounter{probcount}
\newcommand{\exercise}[1]{%
    \stepcounter{probcount}%
    \begin{minipage}[t][110mm][t]{0.44\textwidth}
        \small 
        \vspace{1mm}
        \noindent \textbf{문제 \arabic{probcount}.} #1
    \end{minipage}% 
    }
\setstretch{1.3}

\begin{document}

\noindent
\begin{longtable}{p{0.47\textwidth} | p{0.47\textwidth}}
    
    % [기본 문제 섹션 - 6문항 선별]
    \exercise{다음 정적분의 값을 구하여라.
    \begin{enumerate} [label=(\arabic*), leftmargin=8mm, nosep]
        \item $\displaystyle \int_{0}^{\pi} |\sin x + \cos x| \, dx$
        \item $\displaystyle \int_{0}^{1} \frac{e^x - 1}{e^x + 1} \, dx$
        \item $\displaystyle \int_{1}^{4} e^{\sqrt{x}} \, dx$
    \end{enumerate}}
    & 
    \exercise{자연수 $n$에 대하여 $a_n = \int_{0}^{\frac{\pi}{2}} \sin x \cos x (1-\sin x)^n \, dx$일 때, $\sum_{n=1}^{\infty} a_n$의 값을 구하여라.} \\
    \hline 
    
    \exercise{미분가능한 함수 $f(x)$가 모든 실수 $x$에 대하여 $f(1-x) = 1-f(x)$를 만족시킬 때, 다음을 증명하여라.
\begin{enumerate} [label=(\arabic*), leftmargin=8mm, nosep]
    \item $f'(0) = f'(1)$
    \item $\int_{0}^{1} f(x) \, dx = \frac{1}{2}$
\end{enumerate}} 
    & 
    \exercise{함수 $f(x)$가 모든 실수 $x$에 대하여 $f(x) = f(x+2)$를 만족시키고, $-1 \le x \le 1$일 때 $f(x) = |x|$이다. 이때, $\int_{0}^{2} e^{-x} f(x) \, dx$의 값을 구하여라.} \\
    
    \exercise{다음을 만족시키는 미분가능한 함수 $f(x)$를 구하여라.
\begin{equation*}
    f(x) = x + \int_{0}^{f(0)} e^t f(t) \, dt
\end{equation*}} 
    & 
    \exercise{다음 정적분의 값을 구하여라. (단, $a, b$는 상수이다.)
\begin{enumerate} [label=(\arabic*), leftmargin=8mm, nosep]
    \item $\displaystyle \int_{0}^{1} \frac{1}{(1+x^2)^2} \, dx$
    \item $\displaystyle \int_{0}^{2\pi} |a \sin x + b \cos x| \, dx$ (단, $a, b$는 상수)
    \item $\displaystyle \int_{1}^{e} (1+x)e^x \ln x \, dx$
\end{enumerate}} \\
\hline

    % [실력 문제 섹션 - 5문항 선별]
    \exercise{다음 물음에 답하여라.
\begin{enumerate} [label=(\arabic*), leftmargin=8mm, nosep]
    \item 연속함수 $f(x)$에 대하여 다음이 성립함을 증명하여라.
    \[ \int_{0}^{\frac{\pi}{2}} f(\sin x) \, dx = \int_{0}^{\frac{\pi}{2}} f(\cos x) \, dx \]
    \item $\displaystyle \int_{0}^{\frac{\pi}{2}} \frac{\cos x}{\sin x + \cos x} \, dx$ 의 값을 구하여라.
\end{enumerate}} 
    & 
    \exercise{다음 물음에 답하여라.
\begin{enumerate}[label=(\arabic*), leftmargin=8mm, nosep]
    \item $x = \frac{\pi}{4} - t$로 치환하여 $\displaystyle \int_{0}^{\frac{\pi}{4}} \ln(1 + \tan x) \, dx$의 값을 구하여라.
    \item $\displaystyle \int_{0}^{1} \frac{\ln(1+x)}{1+x^2} \, dx$의 값을 구하여라.
\end{enumerate}} \\
    
    \exercise{다음 함수가 최소가 되는 실수 $t$의 값을 구하여라.
\begin{enumerate}[label=(\arabic*), leftmargin=8mm, nosep]
    \item $f(t) = \int_{1}^{e} (t-x)^2 \ln x \, dx$
    \item $f(t) = \int_{0}^{1} (e^{tx} - 2tx) \, dx$ (단, $t \ge 0$)
\end{enumerate}}
    & 
    \exercise{모든 일차함수 $f(x)$에 대하여 
\[ \int_{0}^{\pi} f(x)(a \sin x + b \cos x + 1) \, dx = 0 \]
이 성립하도록 상수 $a, b$의 값을 정하여라.} \\
\hline

    \exercise{자연수 $n$에 대하여 다음과 같이 정의된 함수 $f_n(x)$가 있다.
\[ f_1(x) = xe^x + \frac{1}{2}, \quad f_{n+1}(x) = xe^x + \frac{1}{2} \int_{0}^{1} f_n(x) \, dx \]
함수 $f_n(x)$에 대하여 $a_{n+1} = \dfrac{1}{2} \int_{0}^{1} f_n(x) \, dx, a_1 = \frac{1}{2}$이라고 할 때,
\begin{enumerate}[label=(\arabic*), leftmargin=8mm, nosep]
    \item $a_{n+1}$을 $a_n$으로 나타내어라.
    \item $\displaystyle \lim_{n \to \infty} a_n$의 값을 구하여라.
\end{enumerate}}
    & 
    \\ % 빈 칸으로 유지하여 균형을 맞춤

\end{longtable}

\end{document}